\clearpage\fancyhead[L]{\small COMET.}\hypertarget{native-390-lower}{}\label{native-390-lower}
\begin{center}\textbf{COMET.}\par\vspace{4pt}\hrule\vspace{4pt}\fontsize{10}{12}\selectfont Göttingische gelehrte Anzeigen. Stück 105. Seite 1041..1043. 1815 Juli 3.\par\vspace{4pt}\hrule\vspace{4pt}\end{center}

\par\indent Im 55. Stück dieser Blätter [S. 382 d. B.] sind die ersten auf der hiesigen Sternwarte
angestellten Beobachtungen des jetzt sichtbaren Cometen mitgetheilt worden: wir können denselben
jetzt noch folgende zwei vom Hrn. Prof. \textsc{Gauss} beifügen, den zufällige Umstände zu einer etwas
langen Unterbrechung genöthigt hatten.

\begin{center}\fontsize{10}{12}\selectfont\setlength{\tabcolsep}{4pt}\renewcommand{\arraystretch}{1.15}\begin{tabular}{l|c|c|c}
1815 & Mittlere Zeit & Gerade Aufst. & Abweichung.\\\hline
April 2 & $9^{\mathrm u}\;11^{\mathrm m}\;55^{\mathrm s}$ & $61^\circ\;2'\;27''$ & $45^\circ\;39'\;57''\;\mathrm{N.}$\\
Juni 12 & $10\;53\;31$ & $165\;3\;3$ & $52\;8\;25$\\
\end{tabular}\end{center}

\Needspace{23\baselineskip}\par\indent Von zahlreichen Beobachtungen, welche von auswärtigen Astronomen mitgetheilt sind, führen wir
hier nur diejenigen an, welche auf der Sternwarte zu Padua von Hrn. \textsc{Santini} angestellt und durch
den Secretär der dortigen Section des K. K. Instituts Hrn. Ritter \textsc{Brera} in einem Schreiben vom 19.
Mai an Hrn Prof. \textsc{Gauss} eingesandt sind.

\SourceLeaf{391}\begin{center}\fontsize{10}{12}\selectfont\setlength{\tabcolsep}{5pt}\renewcommand{\arraystretch}{1.15}\begin{tabular}{l|c|c|c}
1815 & \shortstack{Mittlere Zeit\\in Padua} & \shortstack{Gerade\\Aufsteigung} & \shortstack{Nordliche\\Abweichung}\\\hline
April 24 & $9^{\mathrm u}\;39^{\mathrm m}\;50^{\mathrm s}$ & $80^\circ\;49'\;33''$ & $56^\circ\;3'\;16''$\\
28 & $8\;57\;49$ & $85\;57\;54$ & $57\;37\;53$\\
Mai 1 & $8\;57\;22$ & $90\;16\;13$ & $58\;41\;7$\\
2 & $9\;8\;51$ & $91\;47\;15$ & $59\;0\;32$\\
2 & $9\;40\;22$ & $91\;48\;48$ & $59\;0\;40$\\
6 & $9\;29\;48$ & $98\;13\;9:$ & $\text{wolkig}$\\
8 & $9\;10\;55$ & $101\;40\;4$ & $60\;35\;36$\\
8 & $9\;45\;26$ & $101\;42\;13$ & $60\;37\;0$\\
10 & $9\;9\;11$ & $105\;14\;41$ & $60\;59\;45:$\\
11 & $8\;37\;5$ & $107\;4\;0$ & $61\;6\;4$\\
11 & $9\;22\;30$ & $107\;7\;18$ & $61\;7\;30$\\
12 & $9\;3\;20$ & $108\;57\;37$ & $61\;16\;32:$\\
12 & $9\;48\;15$ & $109\;2\;11$ & $61\;16\;7$\\
17 & $10\;15\;15$ & $118\;47\;45$ & $61\;30\;45$\\
17 & $10\;31\;5$ & $118\;49\;35$ & $61\;30\;47$\\
\end{tabular}\end{center}

\par\indent Diese Beobachtungen wurden mit einem Kreismicrometer am \textsc{Adams}schen Quadranten angestellt.

\par\indent Mehrere Astronomen haben seitdem neue parabolische Elemente bestimmt, welche alle von den
ersten des Hrn. Prof. \textsc{Gauss}, welche a. a. O. mitgetheilt sind, wenig abweichen. So gut aber immer
die frühern Beobachtungen durch alle diese Elemente dargestellt wurden, so hatten diese doch das
Schicksal gemein, dass die spätern Beobachtungen sich schneller und stärker davon entfernten, als
man es sonst bei dergleichen Rechnungen gewohnt ist. Dies deutete schon auf eine merkliche
Verschiedenheit der Bahn von einer Parabel hin. Da der Comet jetzt immer lichtschwächer, und die
Beobachtungen immer schwieriger werden, so hielt Herr Prof. \textsc{Gauss} es für nützlich, zur
Erleichterung der noch zu machenden Beobachtungen noch eine Correction der parabolischen Elemente
vorzunehmen. Er fand aus den Beobachtungen vom 6. März, 25. April (beide von Hrn. Dr. \textsc{Olbers}) und
vom 12. Juni folgende Resultate:

\Needspace{12\baselineskip}\begin{center}\fontsize{10}{12}\selectfont Durchgang durch die Sonnennähe April 25.\par $11^{\mathrm u}\;41^{\mathrm m}\;19^{\mathrm s}$ M. Z. in Göttingen\end{center}
\begin{center}\fontsize{10}{12}\selectfont\setlength{\tabcolsep}{3pt}\renewcommand{\arraystretch}{1.15}\begin{tabular}{lcr}
Länge der Sonnennähe & $.\;.\;.\;.\;.\;.\;.\;.\;.\;.\;.\;.$ & $147^\circ\;35'\;55''$\\
Kleinster Abstand & $.\;.\;.\;.\;.\;.\;.\;.\;.\;.\;.\;.\;.\;.$ & $1.23024$\\
Aufsteigender Knoten & $.\;.\;.\;.\;.\;.\;.\;.\;.\;.\;.\;.\;.$ & $82^\circ\;43'\;6''$\\
Neigung der Bahn & $.\;.\;.\;.\;.\;.\;.\;.\;.\;.\;.\;.\;.$ & $44\;43\;13$\\
\end{tabular}\end{center}\begin{center}Bewegung rechtläufig.\end{center}

\par\indent Diese Elemente stellen die drei Beobachtungen vollständig sehr nahe dar: dessen ungeachtet
entfernen sie sich sehr stark (in den geraden Aufsteigungen \emph{bis auf eilf Minuten}) von den
dazwischen liegenden Beobachtungen. Diese Erscheinung, die oft vorkommen kann, darf nicht
befremden. Es folgt daraus einerseits, dass die Bahn bedeutend von einer Parabel abweicht,
anderseits, dass zur Bestimmung des Kegelschnitts obige drei Beobachtungen nicht geeignet sind,
indem sie allein ihn gewissermassen unbestimmt lassen (Der eigentliche Grund dieses Phänomens
liegt in dem Umstande, dass um die Zeit der mittlern Beobachtung die Richtung der geocentrischen
Bewegung nahe zusammenfiel mit der Richtung eines vom Cometen nach der Sonne gezogenen grössten
Kreises. M. s. die letzten Artikel im ersten Abschnitte des zweiten Buchs der \emph{Theor. Mot.
C. C.}). Es lässt sich daher erwarten, dass folgende kleine nach diesen parabolischen Elementen
berechnete Ephemeride auch sofort merklich von der wirklichen Bewegung abweichen wird, und zwar
wird sie die geraden Aufsteigungen zu klein, die Declinationen zu gross geben; inzwischen wird
dieselbe doch zur blossen leichtern Auffindung genau genug sein:

\SourceLeaf{392}\begin{center}\fontsize{10}{12}\selectfont\setlength{\tabcolsep}{5pt}\renewcommand{\arraystretch}{1.15}\begin{tabular}{l|c|c}
\shortstack{$10^{\mathrm u}\;53^{\mathrm m}$\\in Göttingen} & \shortstack{Gerade\\Aufsteigung} & Abweichung\\\hline
Juni 20 & $174^\circ\;46'$ & $46^\circ\;47'\;\mathrm{N.}$\\
24 & $178\;55$ & $43\;56$\\
28 & $182\;41$ & $41\;3$\\
Juli 2 & $186\;7$ & $38\;8$\\
6 & $189\;16$ & $35\;16$\\
10 & $192\;11$ & $32\;25$\\
14 & $194\;53$ & $29\;40$\\
\end{tabular}\end{center}

\par\indent Die Lichtstärke ist am 20. Juni noch $\frac{3}{10}$, am 14. Juli hingegen nur die Hälfte von
derjenigen, welche der Comet am Tage seiner Entdeckung hatte.

\par\indent Obgleich es zu einer genauen Bestimmung der wahren Bahn des Cometen jetzt noch zu früh ist, so
hielt doch Herr Prof. \textsc{Gauss} es für interessant, eine vorläufige Bestimmung zu machen, die
wenigstens eine ungefähre Vorstellung von der unerwartet grossen Abweichung der Bahn von der
Parabel gibt, und ausserdem zur Discussion und Vorbereitung der Beobachtungen zur schärfern
Bestimmung dienen kann. Diese Rechnung ist, der Absicht gemäss, nur flüchtig gemacht, bloss auf
fünf isolirte Beobachtungen gegründet, und dabei auf Praecession, Aberration, Nutation und
Parallaxe gar keine Rücksicht genommen. Eine ängstliche Berücksichtigung dieser Umstände würde
bei einer solchen ersten Bestimmung verlorne Arbeit gewesen sein.

\Needspace{13\baselineskip}\begin{center}
Vorläufige \emph{elliptische} Elemente des Cometen von 1815.
\end{center}

\begin{center}\fontsize{10}{12}\selectfont Durchgang durch die Sonnennähe April 26 $0^{\mathrm u}\;21^{\mathrm m}\;25^{\mathrm s}$\end{center}
\begin{center}\fontsize{10}{12}\selectfont\setlength{\tabcolsep}{3pt}\renewcommand{\arraystretch}{1.15}\begin{tabular}{lcr}
Länge der Sonnennähe & $.\;.\;.\;.\;.\;.\;.\;.\;.\;.\;.$ & $148^\circ\;58'\;48''$\\
Abstand in der Sonnennähe & $.\;.\;.\;.\;.\;.\;.\;.\;.$ & $1.21349$\\
Aufsteigender Knoten & $.\;.\;.\;.\;.\;.\;.\;.\;.\;.\;.$ & $83\;26\;21$\\
Neigung der Bahn & $.\;.\;.\;.\;.\;.\;.\;.\;.\;.\;.\;.\;.$ & $44\;30\;43$\\
Excentricität & $.\;.\;.\;.\;.\;.\;.\;.\;.\;.\;.\;.\;.\;.$ & $0.933149$\\
\end{tabular}\end{center}

\par\indent Es folgt hieraus noch der Abstand in der Sonnenferne 35.091, Umlaufszeit $77\frac{1}{8}$ Jahre,
nahe eben so gross wie die des \textsc{Halley}schen Cometen. Durch diese Elemente sind alle jene groben
Abweichungen der parabolischen Bahn ganz weggeschafft, und Herr Prof. \textsc{Gauss} glaubt daher, sie
immer schon als eine Annäherung zur Wahrheit betrachten zu dürfen. Sollte vielleicht auch bei
künftiger schärferer Rechnung die Umlaufszeit merklich grösser ausfallen, so wird sie doch
schwerlich über 100 Jahre betragen können, und man kann daher diesen Cometen schon jetzt als eine
der merkwürdigsten Erscheinungen in den Annalen der Astronomie betrachten.
